\section{Bounded-formula Booleanization}
\label{sec:booleanization}

\subsection{The bounded language}

Let atomic predicates be equations or fixed primitive-recursive Boolean tests.
A bounded formula is generated by
\[
\begin{aligned}
\phi ::= {}& P(\vec x)\mid \bot\mid
(\phi\wedge\psi)\mid(\phi\vee\psi)\mid(\phi\to\psi)\\
&\mid \forall y<t(\vec x)\,\phi(\vec x,y)
\mid \exists y<t(\vec x)\,\phi(\vec x,y).
\end{aligned}
\]
The formula is \emph{fixed}: its syntax tree does not arrive as an arbitrary
runtime code.  Parameters and bounds may vary, but the compiler is generated
once for this particular tree.

Use Boolean values \(0,1\), with total operations
\[
  \neg b=1-b,\qquad
  b\wedge c=bc,
  \qquad
  b\vee c=\min(1,b+c),
  \qquad
  b\Rightarrow c=(\neg b)\vee c.
\]
All are primitive recursive.

\subsection{Bounded folds}

Define
\begin{align*}
 \operatorname{All}(0,f)&=1,
 &\operatorname{All}(Sn,f)&=
    \operatorname{All}(n,f)\wedge f(n),\\
 \operatorname{Any}(0,f)&=0,
 &\operatorname{Any}(Sn,f)&=
    \operatorname{Any}(n,f)\vee f(n).
\end{align*}
For each fixed function body \(f\), these are primitive-recursive symbols.
There is no internal higher-type argument: the displayed notation abbreviates
a symbol constructed by primitive recursion from the fixed body.

The least bounded witness is likewise primitive recursive:
\[
 \operatorname{least}(n,f)=
 \begin{cases}
  \min\{i<n:f(i)=1\},&\text{if such }i\text{ exists},\\
  n,&\text{otherwise}.
 \end{cases}
\]
It may be implemented by a scan carrying a ``found'' bit and the current
candidate.

\begin{lemma}[Fold reflection]
For every fixed Boolean term \(f(i,\vec x)\), iPRA proves
\begin{align*}
 \operatorname{All}(n,f)=1
 &\leftrightarrow \forall i<n\,[f(i,\vec x)=1],\\
 \operatorname{Any}(n,f)=1
 &\leftrightarrow \exists i<n\,[f(i,\vec x)=1].
\end{align*}
Moreover, if \(\operatorname{Any}(n,f)=1\), then
\(\operatorname{least}(n,f)<n\) and
\(f(\operatorname{least}(n,f),\vec x)=1\).
\end{lemma}

\begin{proof}
Induct on \(n\).  The zero cases reduce to the defining equations and the
empty bounded quantifiers.  In the successor case, unfold one fold equation.
For \(\operatorname{All}\), Boolean conjunction equals one exactly when both
arguments equal one; apply the induction hypothesis to the prefix and join the
last instance.  The \(\operatorname{Any}\) proof is dual.  For least search,
strengthen the induction predicate to include: either no witness has occurred
and the state equals the current bound, or the stored index is the least
witness seen so far.  This predicate is a fixed Boolean combination of
equalities and inequalities, hence quantifier free after its prefix conditions
are folded.
\end{proof}

\subsection{Structural Booleanization}

Define a characteristic term \(\llbracket\phi\rrbracket\) recursively:
\[
\begin{array}{rcl}
\llbracket P\rrbracket &=& \chi_P,\\
\llbracket\bot\rrbracket&=&0,\\
\llbracket\phi\wedge\psi\rrbracket&=&
 \llbracket\phi\rrbracket\wedge\llbracket\psi\rrbracket,\\
\llbracket\phi\vee\psi\rrbracket&=&
 \llbracket\phi\rrbracket\vee\llbracket\psi\rrbracket,\\
\llbracket\phi\to\psi\rrbracket&=&
 \llbracket\phi\rrbracket\Rightarrow\llbracket\psi\rrbracket,\\
\llbracket\forall y<t\,\phi\rrbracket&=&
 \operatorname{All}(t,\lambda y.\llbracket\phi\rrbracket),\\
\llbracket\exists y<t\,\phi\rrbracket&=&
 \operatorname{Any}(t,\lambda y.\llbracket\phi\rrbracket).
\end{array}
\]
Again, each apparent lambda is eliminated at compile time because the formula
is fixed.

\begin{theorem}[Bounded Booleanization]
\label{thm:booleanization}
For every fixed bounded formula \(\phi(\vec x)\), there is a primitive-
recursive Boolean symbol \(b_\phi\) such that iPRA proves
\[
      \phi(\vec x)\quad\leftrightarrow\quad b_\phi(\vec x)=1.
\]
The equation \(b_\phi(\vec x)=1\) is quantifier free and can be used as an
induction invariant.
\end{theorem}

\begin{proof}
Proceed by structural induction on \(\phi\).  Atomic predicates are represented
by their characteristic functions.  The propositional cases follow from the
finite truth tables, each of which is provable by numerical case analysis on
Boolean inputs.  For bounded universal and existential quantifiers, apply the
fold-reflection lemma to the characteristic function supplied by the
induction hypothesis.  Primitive recursiveness follows by closure under
composition and primitive recursion.  The final equation is atomic.
\end{proof}

\subsection{Intuitionistic implication}

The implication clause deserves comment.  For arbitrary propositions, a
Boolean truth value would classically decide the antecedent.  Here the
antecedent and consequent are bounded formulas with explicit characteristic
functions.  iPRA proves that each is equivalent to a Boolean equation; Boolean
equality is decidable by finite case analysis.  Therefore
\[
 (b_\phi=1\to b_\psi=1)
 \leftrightarrow
 (b_\phi\Rightarrow b_\psi)=1
\]
is intuitionistically derivable.  No global law of excluded middle is used.

\subsection{Why arbitrary formula codes are excluded}

A uniform function taking an arbitrary code of a bounded formula and
returning its characteristic function would be an evaluator for a growing
program language.  A clocked syntactic evaluator is possible, but a total PR
evaluator for all PR programs is not.  The framework therefore uses one of two
modes:
\begin{enumerate}
\item compile a fixed formula at the Lean meta-level into a fixed \(\PR\)
symbol and proof; or
\item evaluate a coded formula for a supplied bounded clock and verify its
trace.
\end{enumerate}
This restriction is not a loss for ordinary theorem formalization, where the
formulas are known in advance.

\subsection{Lean formalization path}

The supplied file \code{BoundedFormulaCompiler.lean} defines a meta-syntax of
fixed bounded formulas and records the intended compilation functions.  A
complete version should proceed in four layers:
\begin{enumerate}
\item add canonical Boolean symbols and prove their equations;
\item define \code{allFold}, \code{anyFold}, and \code{leastFold} as
\(\PR\) terms;
\item prove their object-language reflection theorems by \code{PRA.qfInd};
\item recursively produce a pair consisting of \(b_\phi\) and a
\(\Prf\,\PRAAx\) proof of the equivalence.
\end{enumerate}
The fold layer is the only place requiring induction.  The structural compiler
then uses ordinary propositional proof constructors and substitution.
